\scanpage{204}\mspage{153}
isomorphisme canonique

\diagram{\(\unsure{\mathscr X|\mathcal X}_{/x_0}\simeq\int \begin{tikzcd}[ampersand replacement=\&,row sep=0pt,column sep=1.2em,cramped] \& X_{1/x_1} \\ X_{0/x_0}\arrow[ur]\arrow[dr] \& \\ \& X_{2/x_2} \end{tikzcd},\)}

et la flèche $(\varphi_Q)_{/x_0}$ localisée de $\varphi_Q:\mathscr X\mathbin{\unsure{-|\to}}Y$ en \deleted{\unsure{$y$}} $x_0$ s'identifie à la flèche $\varphi(Q_{/x_0})$ :

\diagram{\(\varphi(Q_{/x_0}):\ \int \begin{tikzcd}[ampersand replacement=\&,row sep=0pt,column sep=1.2em,cramped] \& X_{1/x_1} \\ X_{0/x_0}\arrow[ur]\arrow[dr] \& \\ \& X_{2/x_2} \end{tikzcd}\longrightarrow Y_{/y},\)}

\uline{Démonstration} \inserted{du lemme.} On sait que

\[
\mathscr X_{/x_0}\longrightarrow\mathscr X
\]

est \inserted{$\mathrm{Cat}$-}fibrée \inserted{(scindée)} (à fibres discrètes), et comme

\[
\mathscr X\longrightarrow\Phi
\]

est aussi $\mathrm{Cat}$-fibrée (scindée), il s'ensuit qu'il en est de $=$ de la flèche composée

\[
\mathscr X_{/x_0}\longrightarrow\Phi .
\]

Donc \deleted{il faut} \unsure{,} pour ``calculer'' $\mathscr X_{/x_0}$, en \deleted{\unsure{des}} que catégorie sur $\Phi$, il faut expliciter ses \inserted{trois} fibres, et les deux foncteurs de chgt. de base qui les relient. On trouve aussitôt que ces fibres sont can. isomorphes à

\[
X_{0/x_0},\quad X_{1/x_1},\quad X_{2/x_2},
\]

et on admettra (!) que les foncteurs chgt de
% archive sheet number: 204
